% !TEX root = ../../main_without_quantization.tex
\section{Introduction}

The purpose of pruning is to keep the behavior of a trained neural network while modifying its size by removing some of its parameters or biggest parts. In order to study what is eligible to be removed and how the network will be effected by removing these parts it is first needed to explain what are the objects that pruning can work on: Neurons, weights, activations, layers, feed-forward neural network. The vocabulary is not given here as general knowledge.

We begin this chapter by giving a description of a single artificial neuron. This is necessary for pruning because weights and biases are the smallest parameters which are available for training and activation determines where the information should propagate to after weights and biases have been updated. It naturally leads us to multi-layer perceptrons, loss functions, gradient based methods and backpropagation because we measure the performance of a pruning technique by the amount to which the output behavior, or the loss, changes when removing some part of a network\cite{Goodfellow2016,Rumelhart1986}.


\newpage
\section{Parameterization of Neural Networks}
An artificial neuron is a computing unit that takes an input vector, it applies a linear transformation and finally inputs this result to a non-linear activation function. This form of an artificial neuron enables the artificial neural network to learn non-linear mappings. The linear part is also called an affine transformation, which computes a weighted sum of the input vector and then adds a bias. The non-linear activation function is important to enable the artificial neural network to learn complex patterns, since it is responsible for approximate the non-linearities in the data. Examples for common activation functions are ReLU, sigmoid and tanh.

\subsection{The Neuron as a Linear Model}

The artificial neuron, formally introduced by \cite{McCulloch1943}, serves as the fundamental building block of neural networks. Mathematically, it performs an \textbf{affine transformation} on its inputs followed by a non-linear activation $f$ which is represented in Figure~\ref{fig:neuron_affine_unit} and expressed in Eq.\eqref{eq:neuron_affine_transformation}.

\begin{equation}
\label{eq:neuron_affine_transformation}
z = \vect{w}^T \vect{x} + b = \sum_{i=1}^{n} w_i x_i + b
\end{equation}
\noindent Given an input $\vect{x} \in \reals^n$, weights $\vect{w} \in \reals^n$ scale the importance of each feature, while the bias $b \in \reals$ acts as an offset. Figure~\ref{fig:neuron_affine_unit} shows this affine unit before the activation is applied: the weighted inputs are summed with the bias to produce the pre-activation response $z$.

\begin{figure}[htbp]
    \centering
    \includegraphics[width=0.82\textwidth]{neuron_perceptron_with_activation.pdf}
    \caption{A single artificial neuron with input vector $\vect{x}$, weight vector $\vect{w}$, bias $b$, and pre-activation output $z$. The activation function $f$ is applied to $z$ to produce the final output $y$.}
    \label{fig:neuron_affine_unit}
\end{figure}

The term $\vect{w}^T \vect{x}$ measures how strongly the input matches the learned weights. This affine computation is the basic operation repeated across layers in larger neural networks.
